\begin{abstract}
Mobile GUI agents have grown capable enough to handle multi-step tasks on familiar apps, but their behavior on apps they did not encounter during training remains brittle. We argue that the deployment-relevant generalization gap is not within-template parameter robustness, the regime current benchmarks largely measure, but cross-app and cross-category transfer at a small adaptation budget. We propose \textbf{EvoFSM-RL}, a test-time adaptation framework in which the agent's prior knowledge is a two-layer finite state machine---an app-specific lower layer and a category-generic upper layer, separated by an enforced lint pass---and adaptation runs a joint loop that mutates the upper layer via a frozen reference language model while updating a LoRA adapter via group-relative policy optimization. The same loop produces a shared LoRA on the source pool offline and continues to adapt on each target app's small adaptation set at deployment. We evaluate on a 25-app Android benchmark with three independent splits stacked together: source pool vs.\ target apps, near-transfer vs.\ far-transfer categories, and template-disjoint adaptation vs.\ evaluation. Static category knowledge alone lifts Tier-B success rate by $+9.3$~pp over a zero-shot baseline, online upper-layer evolution adds another $+3.7$~pp, and a Tier-C far-transfer panel acts as a built-in null control for the injection mechanism. The full joint variant completes a four-rung ablation that isolates the contribution of structural priors, online context evolution, and weight co-adaptation.
\end{abstract}
